% Fibonacci-Pell nearest gaps: an all-exponent classification and
% quadratic-unit orbit rigidity
%
% Authors:
%   Dr. Denys Dutykh (Mathematics Department, Khalifa University of Science
%   and Technology, Abu Dhabi, UAE)
%   Prof. Laurent Vuillon (Univ. Savoie Mont Blanc, CNRS, LAMA, Chambery,
%   France)
%
\section{Arithmetic dynamics and complete Pell orbits}
\label{sec:zero-defect}

The nearest-gap equation is not an isolated algebraic coincidence. It is the
exact image of an integral dynamical system under a quadratic-unit
semiconjugacy. This section first proves that structural statement for a
general seed, then specialises it to the two Fibonacci cores and obtains the
complete list in \cref{thm:orbit-intro}.

For an integral seed $z=(u,v)$, put

\begin{equation}
\begin{aligned}
 \mathsf S(u,v)&=(v,u+2v),\\
 \mathcal H(u,v)&=u^2+v^2,\\
 \mathcal B(u,v)&=v^2-u^2+2uv,\\
 \mathcal Q(u,v)&=v^2-2uv-u^2,
\end{aligned}
\label{eq:orbit-forms}
\end{equation}

and define

\begin{equation}
 A_z=(v-u)+u\sqrt2.
 \label{eq:seed-unit-map}
\end{equation}

\begin{proposition}[Seed/unit semiconjugacy]
\label{prop:seed-semiconjugacy}
The map $z\mapsto A_z$ is a bijection

\begin{equation}
 \{(u,v)\in\Z^2:0<u\leqslant v\}
 \longrightarrow
 \{a+b\sqrt2\in\Z[\sqrt2]:a\geqslant0,\ b>0\},
 \label{eq:cone-bijection}
\end{equation}

with inverse $a+b\sqrt2\mapsto(b,a+b)$. For every integral seed and every
$t\geqslant0$,

\begin{equation}
 \norm{A_z}=\mathcal Q(z),
 \qquad
 A_{\mathsf S^tz}=\lambda^tA_z.
 \label{eq:semiconjugacy}
\end{equation}
\end{proposition}

\Cref{fig:seed-semiconjugacy} displays the square that
\cref{eq:semiconjugacy} commutes.

\begin{proof}
The inverse in \cref{eq:cone-bijection} is immediate. Direct calculation
gives

\begin{equation}
 \norm{A_z}=(v-u)^2-2u^2=\mathcal Q(z)
\end{equation}

and

\begin{equation}
 A_{\mathsf Sz}=(u+v)+v\sqrt2
 =\lambda\bigl((v-u)+u\sqrt2\bigr).
\end{equation}

Iteration proves \cref{eq:semiconjugacy}.
\end{proof}

For a positive integral seed, $\mathcal Q(z)=0$ would force
$(v-u)^2=2u^2$, which is impossible. Thus the canonical case
$\mathcal Q(z)^2=1$ and the noncanonical case
$\mathcal Q(z)^2>1$ exhaust all positive seeds.

In coefficient columns this intertwining is an integral conjugacy. Indeed,

\begin{equation}
 M=\begin{pmatrix}0&1\\1&2\end{pmatrix},
 \qquad
 U=\begin{pmatrix}1&2\\1&1\end{pmatrix},
 \qquad
 R=\begin{pmatrix}-1&1\\1&0\end{pmatrix}
 \label{eq:conjugacy-matrices}
\end{equation}

are respectively the seed step, multiplication by $\lambda$, and the map
$(u,v)^{\mathsf T}\mapsto(v-u,u)^{\mathsf T}$. Direct multiplication gives

\begin{equation}
 RM=UR.
 \label{eq:integral-conjugacy}
\end{equation}

Since $R\in\mathrm{GL}_2(\Z)$, no lattice information is lost:
$M^2\in\mathrm{SL}_2(\Z)$ is conjugate to the hyperbolic norm-one unit
$\iota(\lambda^2)$.

\begin{figure}[H]
\centering
\begin{tikzpicture}[
  every node/.style={font=\small,align=center},
  obj/.style={draw,rounded corners,inner sep=6pt,minimum width=35mm},
  arrow/.style={-{Latex[length=2mm]},thick}
]
 \node[obj] (z) at (0,1.7) {$z=(u,v)$\\integral seed};
 \node[obj] (stz) at (7.2,1.7) {$\mathsf S^tz$\\evolved seed};
 \node[obj] (az) at (0,0) {$A_z\in\Z[\sqrt2]$\\quadratic integer};
 \node[obj] (lat) at (7.2,0) {$\lambda^tA_z$\\unit translate};
 \draw[arrow] (z)--node[above] {$\mathsf S^t$} (stz);
 \draw[arrow] (z)--node[left] {$A$} (az);
 \draw[arrow] (stz)--node[right] {$A$} (lat);
 \draw[arrow] (az)--node[below] {$\times\,\lambda^t$} (lat);
\end{tikzpicture}
\caption{The exact seed/unit commutative diagram. The lower dynamics is
multiplication in $\Z[\sqrt2]$, not an asymptotic model of the upper
dynamics. Norm gives the signed covariance $\mathcal Q(\mathsf S^tz)=
(-1)^t\mathcal Q(z)$; equivalently, $\mathcal Q^2$ is invariant.}
\label{fig:seed-semiconjugacy}
\end{figure}

There is a second, symmetric-square description. Direct expansion gives

\begin{equation}
 \mathcal B(z)^2+\mathcal Q(z)^2=2\mathcal H(z)^2,
 \label{eq:lorentz-cone}
\end{equation}

and

\begin{equation}
 \begin{pmatrix}
  \mathcal H(\mathsf S^tz)\\
  \mathcal B(\mathsf S^tz)\\
  \mathcal Q(\mathsf S^tz)
 \end{pmatrix}
 =
 \begin{pmatrix}
  C_{2t}&P_{2t}&0\\
  2P_{2t}&C_{2t}&0\\
  0&0&(-1)^t
 \end{pmatrix}
 \begin{pmatrix}
  \mathcal H(z)\\
  \mathcal B(z)\\
  \mathcal Q(z)
 \end{pmatrix}.
 \label{eq:lorentz-action}
\end{equation}

The integral matrix in \cref{eq:lorentz-action} preserves
$2X^2-Y^2-Z^2$. Thus the quadratic observables of the seed lie on an
integral Lorentz cone, and $\mathsf S$ acts through an arithmetic subgroup of
its orthogonal group. This formulation is useful rather than decorative:
the first two coordinates in \cref{eq:lorentz-action} are exactly the
proportionality mechanism behind the orbit/gap equivalence below.

\begin{theorem}[Anchored orbit/gap equivalence]
\label{thm:general-orbit-gap}
Let $z=(u,v)$ have positive integral coordinates and
$\mathcal Q(z)^2>1$. Let $d$ be positive and $t\geqslant0$, and suppose
either $u\leqslant v$ or $t\geqslant1$. Then

\begin{equation}
 \mathcal H(\mathsf S^tz)=P_{d+2t}
 \label{eq:general-orbit-hit}
\end{equation}

if and only if there is an integer $g>0$ such that

\begin{equation}
 \lambda^{d-1}-A_z^2=g\lambda^{-(2t+1)}.
 \label{eq:general-anchored-gap}
\end{equation}

At a hit, if

\begin{equation}
 D=\mathcal H(z)-P_d,
 \qquad
 E=C_d-\mathcal B(z),
 \label{eq:orbit-DE}
\end{equation}

then

\begin{equation}
 D=gP_{2t},
 \qquad
 E=gC_{2t},
 \qquad
 g=\gcd(D,E).
 \label{eq:orbit-proportionality}
\end{equation}

Moreover, \cref{eq:general-anchored-gap} automatically satisfies

\begin{equation}
 A_z^2<\lambda^{d-1}<2A_z^2,
 \label{eq:general-orbit-window}
\end{equation}

and its coefficient content and norm are respectively $g$ and $-g^2$.
\end{theorem}

\begin{proof}
Put $w=\mathsf S^tz$ and $\ell=d+2t$. The hypotheses ensure that $w$ is
positive and ordered. At a hit, \cref{eq:lorentz-cone} and
$C_\ell^2-2P_\ell^2=(-1)^d$ give

\begin{equation}
 \mathcal B(w)^2
 =C_\ell^2-\bigl(\mathcal Q(z)^2+(-1)^d\bigr)<C_\ell^2.
 \label{eq:B-below-C}
\end{equation}

Because $w$ is positive and ordered, $0<\mathcal B(w)<C_\ell$. For
$t\geqslant1$, comparison of the first coordinates in
\cref{eq:lorentz-action} with the Pell addition formula gives

\begin{equation}
 C_{2t}D=P_{2t}E.
\end{equation}

The second-coordinate difference is
$\mathcal B(w)-C_\ell=-E/C_{2t}$, so $D,E>0$. Since
$\gcd(C_{2t},P_{2t})=1$, this proves
\cref{eq:orbit-proportionality}. At $t=0$, the hit gives $D=0$ and
\cref{eq:B-below-C} gives $E>0$; the same formula holds with
$P_0=0$ and $C_0=1$.

The parity-free identity

\begin{equation}
 \lambda A_z^2=\mathcal B(z)+\mathcal H(z)\sqrt2
 \label{eq:key-seed-identity}
\end{equation}

turns \cref{eq:orbit-DE} into

\begin{equation}
 \lambda^d-\lambda A_z^2=E-D\sqrt2.
 \label{eq:gap-from-DE}
\end{equation}

At a hit, \cref{eq:orbit-proportionality} makes the right-hand side
$g\lambda^{-2t}$, and division by $\lambda$ proves
\cref{eq:general-anchored-gap}. Conversely,
\cref{eq:general-anchored-gap,eq:gap-from-DE} recover
\cref{eq:orbit-proportionality} coefficientwise; the first row of
\cref{eq:lorentz-action} then recovers the hit.

It remains to prove the window. Multiplying
\cref{eq:general-anchored-gap} by $\lambda^{2t}$ and using
\cref{eq:semiconjugacy} gives

\begin{equation}
 \lambda^{\ell-1}-A_w^2=g\lambda^{-1}.
 \label{eq:evolved-gap}
\end{equation}

At the hit, $g=C_\ell-\mathcal B(w)$. Combining this with
\cref{eq:key-seed-identity} at $w$ gives

\begin{equation}
 \lambda A_w^2-g
 =2\mathcal B(w)+(-1)^{\ell+1}\lambda^{-\ell}>0.
\end{equation}

For even $\ell$, positivity follows from
$2\mathcal B(w)\geqslant2>\lambda^{-\ell}$; for odd $\ell$ it is immediate.
Hence $0<g\lambda^{-(2t+1)}<A_z^2$, which is
\cref{eq:general-orbit-window}. Finally,
$\lambda^{-(2t+1)}$ is a coefficient-primitive norm-$-1$ unit, proving the
content and norm assertions.
\end{proof}

\begin{remark}[Canonical seeds]
\label{rem:canonical-seeds}
If $0<u\leqslant v$ and $\mathcal Q(z)^2=1$, then
\cref{prop:seed-semiconjugacy} and the unit group of $\Z[\sqrt2]$ give a
unique $r\geqslant1$ with

\begin{equation}
 z=(P_r,P_{r+1}),
 \qquad
 A_z=\lambda^r.
\end{equation}

Consequently

\begin{equation}
 \mathcal H(\mathsf S^tz)=P_{d+2t}
 \quad\Longleftrightarrow\quad
 d=2r+1,
\end{equation}

and this equality persists for every $t\geqslant0$. The corresponding gap
is zero, so these seeds are correctly separated from
\cref{thm:general-orbit-gap}.
\end{remark}

We now specialise to

\begin{equation}
 z_{q,\sigma}=(F_{q+1},F_{q+1}+\sigma F_q).
 \label{eq:fibonacci-cores}
\end{equation}

For both signs,

\begin{equation}
 A_{z_{q,\sigma}}=D_{q,\sigma},
 \qquad
 \mathcal Q(z_{q,\sigma})=F_q^2-2F_{q+1}^2.
 \label{eq:fibonacci-core-map}
\end{equation}

\begin{theorem}[Complete all-anchor Fibonacci-core orbits]
\label{thm:complete-fibonacci-orbits}
For every $q\geqslant1$, $\sigma\in\{1,-1\}$, positive $d$, and
$t\geqslant0$, the hit

\begin{equation}
 \mathcal H(\mathsf S^tz_{q,\sigma})=P_{d+2t}
 \label{eq:fibonacci-orbit-hit}
\end{equation}

occurs exactly in the four rows of \cref{eq:intro-orbit-list}.
\end{theorem}

\begin{proof}
First let $q\geqslant2$. Then

\begin{equation}
 \mathcal Q(z_{q,\sigma})^2-1
 =F_{q-1}F_{q+1}F_{q+2}F_{q+4}>0.
 \label{eq:core-noncanonical}
\end{equation}

The strict window in \cref{thm:general-orbit-gap} first forces
$d\geqslant2$: if $d=1$, it would give
$A_{z_{q,\sigma}}^2<1$, whereas $A_{z_{q,\sigma}}>1$ for
$q\geqslant2$. Hence an odd anchor here has $d\geqslant3$, so $d-1$ is a
positive even exponent.

For the plus sign the seed is ordered, so
\cref{thm:general-orbit-gap} applies at every time. If $d$ is odd, the gap
exponent $d-1$ is even. The only plus target in
\cref{thm:classification-intro} has normalised direction $\lambda$, whereas
the orbit bridge requires $\lambda^{-(2t+1)}$. Hence there is no plus hit.
If $d$ is even, the gap exponent is odd, which is impossible by
\cref{thm:classification-intro}. Thus the plus core has no hit for
$q\geqslant2$.

For the minus sign and $t\geqslant1$, the evolved seed is ordered, so the
same bridge applies. At an odd anchor, the only minus target has normalised
direction $\lambda^{-1}$; it would force $t=0$, a contradiction. At an even
anchor, the exponent is odd and no target exists. Thus the minus core has no
positive-time hit.

It remains to classify direct minus hits. The identity

\begin{equation}
 \mathcal H(z_{q,-1})
 =F_{q+1}^2+F_{q-1}^2
 =3F_q^2+2(-1)^q
 \label{eq:minus-direct-height}
\end{equation}

and the complete Pell-value classification used in
\cref{prop:negative-exponent} give exactly

\begin{equation}
 (q,d)=(2,3),(4,5)
 \label{eq:minus-direct-hits}
\end{equation}

for $q\geqslant2$. These are the last two rows of
\cref{eq:intro-orbit-list}.

Finally, $z_{1,1}=(1,2)=(P_1,P_2)$ is canonical. By
\cref{rem:canonical-seeds}, it gives $d=3$ and every $t\geqslant0$.
For the minus sign, $z_{1,-1}=(1,0)$ gives the direct hit
$\mathcal H(1,0)=P_1$. If $t\geqslant1$, then

\begin{equation}
 \mathsf S^t(1,0)=(P_{t-1},P_t),
 \qquad
 \mathcal H(\mathsf S^t(1,0))=P_{2t-1}<P_{d+2t},
\end{equation}

so there are no further hits. This proves the exhaustive list.
\end{proof}

\begin{remark}[Exact scope]
The theorem classifies both signs, every positive anchor, and every
nonnegative time for the displayed Fibonacci cores. It does not classify
negative time or all noncanonical positive seeds with
$\mathcal Q(z)^2>1$. Most importantly, no result here shows that a
hypothetical Markoff collision must produce one of these cores. The
semiconjugacy closes the downstream arithmetic problem; it does not supply
the missing upstream Markoff reduction.
\end{remark}

The next section changes perspective. Instead of classifying the global
equation, it asks how much of its divisibility can recur locally, and thereby
identifies the precise limits of $p$-adic information.
